% 图契约：欧姆刻度来自电流的反比例关系；图中的端口、几何、箭头和方程对应本题。
% 版面：沿用前三模块的几何图+推导分区；目标印宽17cm，简图至中等密度。
% 布局：左侧真实结构，右侧关键关系；真实电路边不作为架构流程箭头处理。
\documentclass[tikz,border=7pt]{standalone}
\usepackage{fontspec,xeCJK,amsmath,amssymb}
\IfFontExistsTF{Times New Roman}{\setmainfont{Times New Roman}}{\setmainfont{TeX Gyre Termes}}
\IfFontExistsTF{Heiti SC}{\setCJKmainfont[AutoFakeBold=1.5]{Heiti SC}}{\setCJKmainfont[AutoFakeBold=1.5]{FandolHei-Regular}}
\usetikzlibrary{arrows.meta,bending,calc,patterns,positioning,decorations.pathmorphing}
\definecolor{ink}{HTML}{243745}
\definecolor{blue}{HTML}{275DAB}
\definecolor{green}{HTML}{177C63}
\definecolor{red}{HTML}{BD4935}
\definecolor{orange}{HTML}{AA6C16}
\definecolor{purple}{HTML}{715393}
\tikzset{every node/.style={font=\fontsize{9.3}{12.5}\selectfont,text=ink},
arrow/.style={-{Stealth[length=5pt 1.5,width'=0pt 3]},shorten >=2pt,shorten <=1pt,line width=1pt},
arrow short/.style={-{Stealth[length=3pt 1.0,width'=0pt 2.5]},shorten >=1pt,shorten <=0pt,line width=.8pt},
axis/.style={arrow,draw=ink!60,line width=.6pt},
vec/.style={arrow,draw=blue},
wire/.style={draw=ink!75,line width=.8pt},
guide/.style={dashed,draw=ink!45,line width=.5pt},
card/.style={rounded corners=2pt,draw=blue!65,fill=blue!4,align=center,text width=6.5cm,inner xsep=9pt,inner ysep=7pt},
title/.style={font=\fontsize{10.5}{14}\selectfont\bfseries,text width=15cm,align=center},
label/.style={font=\fontsize{8.4}{11}\selectfont,align=center},
note/.style={label,text width=6.1cm}}
\begin{document}
\begin{tikzpicture}[x=1cm,y=1cm]
\path[use as bounding box] (-8,-3.5) rectangle(8,3.65);
\node[title] at (0,3.25) {M4-037\quad 欧姆刻度来自电流的反比例关系};
\begin{scope}[shift={(-4.05,0)}]
\draw[wire,draw=ink](-2,1.2)--(-0.475,1.2) (0.07500000000000007,1.2)--(1.6,1.2);\begin{scope}[shift={(-0.19999999999999996,1.2)},rotate=0.0]\draw[draw=ink,fill=white,line width=.8pt](-0.275,-.085)rectangle(0.275,.085);\end{scope}
\node[label] at(-0.19999999999999996,1.48){$R_0+R_x$};
\draw[wire](1.6,1.2)--(2,1.2)--(2,-1.2)--(-2,-1.2)--(-2,-.4);\draw[wire](-2,.4)--(-2,1.2);\draw[wire](-2,0)circle(.4);\node at(-2,.17){$+$};\node at(-2,-.17){$-$};\node[left=5pt]at(-2,0){$\mathcal E$};\draw[vec,red](.8,-1.8)--(-.8,-1.8)node[midway,below]{$I$};\node[label]at(0,0){串联式欧姆表};\end{scope}
\node[card,draw=blue!65,fill=blue!4] at (3.9,1.8) {$I_g=\mathcal E/R_0$};
\node[card,draw=green!65,fill=green!4] at (3.9,0) {$I/I_g=R_0/(R_0+R_x)$};
\node[card,draw=orange!65,fill=orange!4] at (3.9,-1.8) {$I=I_g/2\ \Rightarrow\ R_x=R_0$};
\end{tikzpicture}\end{document}